% 计算机学报模板 - XeLaTeX 适配版
% 修改说明：
% 1. 改用 xeCJK 管理中文，删除 CJKutf8 及所有 CJK 环境
% 2. 删除所有 strip 环境（解决 \end{strip} 报错）
% 3. 设置中文字体：宋体、黑体、仿宋、楷体（请根据系统字体调整）
% 4. 保留原模板样式与自定义命令
% 5. 编译方式：XeLaTeX

\documentclass[10.5pt,compsoc]{CjC}
\usepackage{xeCJK}                 % 中文支持

% ------------------ 中文字体设置（请根据您的系统修改）------------------
% Windows 典型字体名：SimSun, SimHei, FangSong, KaiTi
% macOS 典型字体名：Songti SC, STHeiti, STFangsong, STKaiti
% 如使用上传的字体文件，请用具体文件名，例如：
% \setCJKmainfont{simsun.ttc}[AutoFakeBold=2.5]
\setCJKmainfont{simsun.ttc}[AutoFakeBold=2.5]  % 宋体（正文），启用伪粗体
\setCJKsansfont{simhei.ttf}[AutoFakeBold=2.5]  % 黑体（标题等）
\setCJKmonofont{simfang.ttf}                   % 仿宋（作者、摘要等）
\setCJKfamilyfont{zhkai}{simkai.ttf}           % 楷体（致谢等）
\newcommand{\kai}{\CJKfamily{zhkai}}           % 方便调用楷体
% ----------------------------------------------------------------------

% 原模板中的 zhwinfonts.tex 不再需要（已被 xeCJK 替代），可注释或删除
% \input{zhwinfonts}

\usepackage{graphicx}
\usepackage{footmisc}
\usepackage{subfigure}
\usepackage{url}
\usepackage{multirow}
\usepackage[noadjust]{cite}
\usepackage{amsmath,amsthm}
\usepackage{amssymb,amsfonts}
\usepackage{booktabs}
\usepackage{color}
\usepackage{ccaption}
\usepackage{float}
\usepackage{fancyhdr}
\usepackage{caption}
\usepackage{xcolor,stfloats}
\usepackage{comment}
\usepackage{cuted}       % 提供 cuted 环境（但不再使用 strip）
\usepackage{captionhack}
\usepackage{epstopdf}
% 注意：不再需要 inputenc 和 CJKutf8

\setcounter{page}{1}
\graphicspath{{figures/}}

% ------------------ 原模板中的页眉定义（去除 CJK 环境）------------------
\headevenname{\mbox{\quad} \hfill  \mbox{\zihao{-5}{计\quad \quad 算\quad \quad 机\quad \quad 学\quad \quad 报}} \hspace {50mm} \mbox{2019 年}}%
\headoddname{? 期 \hfill {{AUTHOR}}：{{TITLE}}}%

% 脚注编号样式（原模板使用符号，改为数字便于每页重置，但每页重置需额外包）
\renewcommand{\thefootnote}{\arabic{footnote}}
\setcounter{footnote}{0}
\renewcommand\footnotelayout{\zihao{5-}}

\newtheoremstyle{mystyle}{0pt}{0pt}{\normalfont}{1em}{\bf}{}{1em}{}
\theoremstyle{mystyle}
\renewcommand\figurename{figure~}
\renewcommand{\thesubfigure}{(\alph{subfigure})}
\newcommand{\upcite}[1]{\textsuperscript{\cite{#1}}}
\renewcommand{\labelenumi}{(\arabic{enumi})}
\newcommand{\tabincell}[2]{\begin{tabular}{@{}#1@{}}#2\end{tabular}}
\newcommand{\abc}{\color{white}\vrule width 2pt}
\makeatletter
\renewcommand{\@biblabel}[1]{[#1]\hfill}
\makeatother
\setlength\parindent{2em}

% 字号命令（原模板已有，保留）
\makeatletter
\newcommand\mysmall{\@setfontsize\mysmall{7}{9.5}}
\newenvironment{tablehere}{\def\@captype{table}}{}
\let\temp\footnote
\renewcommand \footnote[1]{\temp{\zihao{-5}#1}}
\makeatother

\begin{document}

\hyphenpenalty=50000
\thispagestyle{plain}%
\thispagestyle{empty}%
\pagestyle{CjCheadings}

% ------------------ 文章头部（原 table* 内容，去除 CJK 环境）------------------
\begin{table*}[!t]
	\vspace {-13mm}
	\begin{tabular}{p{168mm}}
		\zihao{5-}
		第??卷\quad 第?期 \hfill 计\quad 算\quad 机\quad 学\quad 报\hfill Vol. ??  No. ?\\
		\zihao{5-}
		{{DATE}} \hfill CHINESE JOURNAL OF COMPUTERS \hfill {{DATE}}\\
		\hline\\[-4.5mm]
		\hline
	\end{tabular}
	
	\centering
	\vspace {11mm}
	{\zihao{2} {{TITLE}} }
	\vskip 5mm
	
	{\zihao{3} {{AUTHOR}}}
	
	\vspace {5mm}
	\zihao{6}{$^{1)}$(单位全名 部门(系)全名, 市(或直辖市) 国家名 邮政编码) *字体为6号宋体*单位}
	
	\zihao{6}{$^{2)}$(单位全名 部门(系)全名, 市(或直辖市) 国家名 邮政编码)*中英文单位名称、作者姓名须一致*}
	
	\zihao{6}{$^{3)}$(单位全名 部门(系)全名, 市(或直辖市) 国家名 邮政编码)}
	
	\zihao{6}{论文定稿后，作者署名、单位无特殊情况不能变更。若变更，须提交签章申请，国家名为中国可以不写，省会城市不写省的名称，其他国家必须写国家名。}
	
	\vskip 5mm
	{\centering
		\begin{tabular}{p{160mm}}
			\zihao{5-}{
				\setlength{\baselineskip}{16pt}\selectfont{
					\noindent 摘\quad 要\quad  {{ABSTRACT}}
					\par}}\\[2mm]
			
			\zihao{5-}{\noindent
				关键词 \quad {{KEYWORDS}}
			}\\[2mm]
			\zihao{5-}{中图法分类号	TP \quad \quad \quad     DOI号:*投稿时不提供DOI号}
	\end{tabular}}
	
	\vskip 7mm
	
	\begin{center}
		\zihao{3}{ {{TITLE}}}\\
		\vspace {5mm}
		\zihao{5}{ {{AUTHOR}}}\\
		\vspace {2mm}
		\zihao{6}{$^{1)}$(Department of ****, University, City ZipCode, China) *字体为6号Times new Roman* Depart.Correspond}
		
		\zihao{6}{$^{2)}$(Department of ****, University, City ZipCode)*中国不写国家名*}
		
		\zihao{6}{$^{3)}$(Department of ****, University, City ZipCode, country)*外国写国家名*}
	\end{center}
	
	\begin{tabular}{p{160mm}}
		\zihao{5}{
			\setlength{\baselineskip}{18pt}\selectfont{
				{\bf Abstract}\quad {{ABSTRACT}}
				\par}}\\
		
		\setlength{\baselineskip}{18pt}\selectfont{
			\zihao{5}{\noindent {{ABSTRACT}}
				
				\vspace {5mm}
				{\bf Keywords}\quad {{KEYWORDS}}\par}
		}
	\end{tabular}
	
	\setlength{\tabcolsep}{2pt}
	\begin{tabular}{p{0.05cm}p{16.15cm}}
		\multicolumn{2}{l}{\rule[4mm]{40mm}{0.1mm}}\\[-3mm]
		&
		收稿日期：{{DATE}}；最终修改稿收到日期：\quad \quad -\quad -\quad .*投稿时不填写此项*. 本课题得到… …基金中文完整名称(No.项目号)、… …基金中文完整名称(No.项目号)、… … 基金中文完整名称(No.项目号)资助.{{AUTHOR}}(通信作者)，性别，xxxx年生，学位(或目前学历)，职称，是/否计算机学会(CCF)会员（提供会员号),主要研究领域为*****、****.E-mail: **************.(给出的电子邮件地址应不会因出国、毕业、更换工作单位等原因而变动。请给出所有作者的电子邮件)
		第1作者手机号码(投稿时必须提供，以便紧急联系，发表时会删除): … …, E-mail: … …*此部分6号宋体*
	\end{tabular}
\end{table*}
\clearpage

% 正文开始
\linespread{1.15}
\zihao{5}
\vskip 1mm
{{BODY}}

\vspace {3mm}
\zihao{5}{
	\noindent 致\quad 谢\quad *致谢内容.* 致谢}

\vspace {5mm}
\centerline
{\zihao{5}
	参~考~文~献}

\begin{thebibliography}{99}
	\zihao{5-} \addtolength{\itemsep}{-1em}
	\vspace {1.5mm}
	{{REFERENCES}}
\end{thebibliography}

\noindent {\zihao{5}\bf{附录X}.}

{\zihao{5-}\setlength\parindent{2em}
	{{APPENDIX}}}

\begin{biography}[yourphotofilename.jpg]
	\noindent
	\textbf{First A. Author}\ \ *计算机学报第1作者提供照片电子图片，尺寸为1寸。英文作者介绍内容包括：出生年,学位(或目前学历),职称,主要研究领域(\textbf{与中文作者介绍中的研究方向一致}).*
	*字体为小5号Times New Roman*
\end{biography}

\begin{biography}[yourphotofilename.jpg]
	\noindent
	\textbf{Second B. Author} *英文作者介绍内容包括：出生年,学位(或目前学历),职称,主要研究领域(\textbf{与中文作者介绍中的研究方向一致})。*
	*字体为小5号Times New Roman*
\end{biography}

\zihao{5}
\noindent \textbf{Background}

\zihao{5-}{
	\setlength\parindent{2em}
	*论文背景介绍为\textbf{英文}，字体为小5号Times New Roman体*
	
	论文后面为400单词左右的英文背景介绍。介绍的内容包括：
	
	本文研究的问题属于哪一个领域的什么问题。该类问题目前国际上解决到什么程度。
	
	本文将问题解决到什么程度。
	
	课题所属的项目。
	
	项目的意义。
	
	本研究群体以往在这个方向上的研究成果。
	
	本文的成果是解决大课题中的哪一部分，如果涉及863$\backslash$973以及其项目、基金、研究计划，注意这些项目的英文名称应书写正确。
}

\end{document}
